\begin{textsample}{POS Dim 1 – 2007 – Score 34.00 – 2007/t000298.txt}  \label{ex:f1_pos_020}
As an architect you design for \textbf{the} present, with an awareness of \textbf{the} past, for a future which is essentially unknown. \textbf{The} green agenda is probably \textbf{the} most important agenda and issue of \textbf{the} day.
And I ' d like to share some experience over \textbf{the} last 40 years — we celebrate our fortieth anniversary this year — and to explore and to touch on some observations about \textbf{the} nature of \textbf{sustainability}.
How far you can anticipate, what follows from it, what are \textbf{the} threats, what are \textbf{the} possibilities, \textbf{the} challenges, \textbf{the} opportunities?
I think that — I ' ve said in \textbf{the} past, many, many years ago, before anybody even invented \textbf{the} concept of a green agenda, that it wasn ' t about fashion — it was about survival.
But what I never said, and what I ' m really going to make \textbf{the} point is, that really, green is cool.
I mean, all \textbf{the} projects which have, in some way, been inspired by that agenda are about a celebratory lifestyle, in a way celebrating \textbf{the} places and \textbf{the} spaces which determine \textbf{the} quality of life.
I rarely actually quote anything, so I ' m going to try and find a piece of paper if I can, [ in ] which somebody, at \textbf{the} end of last year, ventured \textbf{the} thought about what for that individual, as a kind of important observer, analyst, writer — a guy called Thomas Friedman, who wrote in \textbf{the} Herald Tribune, about 2006.
He said,"I think \textbf{the} most important thing to happen in 2006 was that living and thinking green hit Main Street.
We reached a tipping point this year where living, acting, designing, investing and manufacturing green came to be understood by a critical mass of citizens, entrepreneurs and officials as \textbf{the} most patriotic, capitalistic, geo-political and competitive thing they could do.
Hence my motto : green is \textbf{the} new red, white and \textbf{blue}."And I asked myself, in a way, looking back,"When did that kind of awareness of \textbf{the} \textbf{planet} and its fragility first appear?"And I think it was July 20, 1969, when, for \textbf{the} first time, man could look back at \textbf{planet} Earth.
And, in a way, it was Buckminster Fuller who coined that phrase.
And before \textbf{the} kind of collapse of \textbf{the} communist system, I was privileged to meet a lot of cosmonauts in Space City and other places in Russia.
And interestingly, as I think back, they were \textbf{the} first true environmentalists.
They were filled with a kind of pioneering passion, fired about \textbf{the} problems of \textbf{the} Aral Sea.
And at that period it was — in a way, a number of things were happening.
Buckminster Fuller was \textbf{the} kind of green guru — again, a word that had not been coined.
He was a design scientist, if you like, a poet, but he foresaw all \textbf{the} things that are happening now.
It ' s another subject.
It ' s another conversation.
You can go back to his writings : it ' s quite extraordinary.
It was at that time, with an awareness fired by Bucky ' s prophecies, his concerns as a citizen, as a kind of citizen of \textbf{the} \textbf{planet}, that influenced my thinking and what we were doing at that time.
And it ' s a number of projects.
I select this one because it was 1973, and it was a master plan for one of \textbf{the} Canary Islands.
And this probably coincided with \textbf{the} time when you had \textbf{the} \textbf{planet} Earth ' s sourcebook, and you had \textbf{the} hippie movement.
And there are some of those qualities in this drawing, which seeks to \textbf{sum} up \textbf{the} recommendations.
And all \textbf{the} components are there which are now in common parlance, in our vocabulary, you know, 30-odd years later : wind energy, recycling, biomass, solar cells.
And in parallel at that time, there was a very kind of exclusive design club.
People who were really design conscious were inspired by \textbf{the} work of Dieter Rams, and \textbf{the} \textbf{objects} that he would create for \textbf{the} company called Braun.
This is going back \textbf{the} mid- ' 50s, ' 60s.
And despite Bucky ' s prophecies that everything would be miniaturized and technology would make an incredible style — access to comfort, to amenities — it was very, very difficult to imagine that everything that we see in this image, would be very, very stylishly packaged.
And that, and more besides, would be in \textbf{the} palm of your hand.
And I think that that digital revolution now is coming to \textbf{the} point where, as \textbf{the} virtual world, which brings so many people together here, finally connects with \textbf{the} physical world, there is \textbf{the} reality that that has become humanized, so that digital world has all \textbf{the} friendliness, all \textbf{the} immediacy, \textbf{the} orientation of \textbf{the} analog world.
Probably \textbf{summed} up in a way by \textbf{the} stylish or alternative available here, as we generously had gifted at lunchtime, \textbf{the} [ \textbf{unclear} ], which is a further kind of development — and again, inspired by \textbf{the} incredible sort of sensual feel.
A very, very beautiful \textbf{object}.
So, something which in [ \textbf{the} ] ' 50s, ' 60s was very exclusive has now become, interestingly, quite inclusive.
And \textbf{the} reference to \textbf{the} iPod as iconic, and in a way evocative of performance, delivery — quite interesting that [ in ] \textbf{the} beginning of \textbf{the} year 2007, \textbf{the} Financial Times commented that \textbf{the} Detroit companies envy \textbf{the} halo effect that Toyota has gained from \textbf{the} Prius as \textbf{the} hybrid, energy-conscious vehicle, which rivals \textbf{the} iPod as an iconic product.
And I think it ' s very tempting to, in a way, seduce ourselves — as architects, or anybody involved with \textbf{the} design process — that \textbf{the} answer to our problems lies with buildings.
Buildings are important, but they ' re only a component of a much bigger picture.
In other words, as I might seek to demonstrate, if you could achieve \textbf{the} impossible, \textbf{the} equivalent of perpetual motion, you could design a carbon-free house, for example.
That would be \textbf{the} answer.
Unfortunately, it ' s not \textbf{the} answer.
It ' s only \textbf{the} beginning of \textbf{the} problem.
You cannot separate \textbf{the} buildings out from \textbf{the} infrastructure of cites and \textbf{the} mobility of transit.
For example, if, in that Bucky-inspired phrase, we draw back and we look at \textbf{planet} Earth, and we take a kind of typical, industrialized society, then \textbf{the} energy consumed would be split between \textbf{the} buildings, 44 percent, transport, 34 percent, and industry.
But again, that only shows part of \textbf{the} picture.
If you looked at \textbf{the} buildings together with \textbf{the} associated transport, in other words, \textbf{the} transport of people, which is 26 percent, then 70 percent of \textbf{the} energy \textbf{consumption} is influenced by \textbf{the} way that our cites and infrastructure work together.
So \textbf{the} problems of \textbf{sustainability} cannot be separated from \textbf{the} nature of \textbf{the} cities, of which \textbf{the} buildings are a part.
For example, if you take, and you make a comparison between a recent kind of city, what I ' ll call, simplistically, a North American city — and Detroit is not a bad example, it is very car dependent. \textbf{The} city goes out in annular rings, consuming more and more green space, and more and more roads, and more and more energy in \textbf{the} transport of people between \textbf{the} city center — which again, \textbf{the} city center, as it becomes deprived of \textbf{the} living and just becomes commercial, again becomes dead.
If you compared Detroit with a city of a Northern European example — and Munich is not a bad example of that, with \textbf{the} greater dependence on walking and cycling — then a city which is really only twice as dense, is only using one-tenth of \textbf{the} energy.
In other words, you take these comparable examples and \textbf{the} energy leap is enormous.
So basically, if you wanted to generalize, you can demonstrate that as \textbf{the} density increases along \textbf{the} bottom there, that \textbf{the} energy consumed reduces dramatically.
Of course you can ' t separate this out from issues like social diversity, mass transit, \textbf{the} ability to be able to walk a convenient distance, \textbf{the} quality of civic spaces.
But again, you can see Detroit, in \textbf{yellow} at \textbf{the} top, extraordinary \textbf{consumption}, down below Copenhagen.
And Copenhagen, although it ' s a dense city, is not dense compared with \textbf{the} really dense cities.
In \textbf{the} year 2000, a rather interesting thing happened.
You had for \textbf{the} first time mega-cities, [ of ] 5 million or more, which were \textbf{occurring} in \textbf{the} developing world.
And now, out of typically 46 cities, 33 of those mega-cities are in \textbf{the} developing world.
So you have to ask yourself — \textbf{the} environmental impact of, for example, China or India.
If you take China, and you just take Beijing, you can see on that traffic system, and \textbf{the} pollution associated with \textbf{the} \textbf{consumption} of energy as \textbf{the} cars expand at \textbf{the} price of \textbf{the} bicycles.
In other words, if you put onto \textbf{the} roads, as is currently happening, 1, 000 new cars every day — statistically, it ' s \textbf{the} biggest booming auto market in \textbf{the} world — and \textbf{the} half a billion bicycles serving one and a third billion people are reducing.
And that urbanization is extraordinary, accelerated pace.
So, if we think of \textbf{the} transition in our society of \textbf{the} movement from \textbf{the} land to \textbf{the} cities, which took 200 years, then that same process is happening in 20 years.
In other words, it is accelerating by a factor of 10.
And quite interestingly, over something like a 60-year period, we ' re seeing \textbf{the} doubling in life expectancy, over that period where \textbf{the} urbanization has trebled.
If I pull back from that global picture, and I look at \textbf{the} implication over a similar period of time in terms of \textbf{the} technology — which, as a tool, is a tool for designers, and I cite our own experience as a company, and I just illustrate that by a small \textbf{selection} of projects — then how do you measure that change of technology?
How does it affect \textbf{the} design of buildings?
And particularly, how can it lead to \textbf{the} creation of buildings which consume less energy, create less pollution and are more socially responsible?
That story, in terms of buildings, started in \textbf{the} late ' 60s, early ' 70s. \textbf{The} one example I take is a corporate headquarters for a company called Willis and Faber, in a small market town in \textbf{the} northeast of England, commuting distance with London.
And here, \textbf{the} first thing you can see is that this building, \textbf{the} roof is a very warm kind of overcoat blanket, a kind of insulating garden, which is also about \textbf{the} celebration of public space.
In other words, for this community, they have this garden in \textbf{the} sky.
So \textbf{the} humanistic ideal is very, very strong in all this work, encapsulated perhaps by one of my early sketches here, where you can see greenery, you can see sunlight, you have a connection with nature.
And nature is part of \textbf{the} generator, \textbf{the} driver for this building.
And symbolically, \textbf{the} colors of \textbf{the} interior are green and \textbf{yellow}.
It has facilities like swimming pools, it has flextime, it has a social heart, a space, you have contact with nature.
Now this was 1973.
In 2001, this building received an award.
And \textbf{the} award was about a celebration for a building which had been in use over a long period of time.
And \textbf{the} people who ' d created it came back : \textbf{the} project managers, \textbf{the} company chairmen then.
And they were saying, you know,"\textbf{The} architects, Norman was always going on about designing for \textbf{the} future, and you know, it didn ' t seem to cost us any more.
So we humored him, we kept him happy."\textbf{The} image at \textbf{the} top, what it doesn ' t — if you look at it in detail, really what it is saying is you can wire this building.
This building was wired for change.
So, in 1975, \textbf{the} image there is of typewriters.
And when \textbf{the} photograph was taken, it ' s word processors.
And what they were saying on this occasion was that our competitors had to build new buildings for \textbf{the} new technology.
We were fortunate, because in a way our building was future-proofed.
It anticipated change, even though those changes were not known.
Round about that design period leading up to this building, I did a sketch, which we pulled out of \textbf{the} archive recently.
And I was saying, and I wrote,"But we don ' t have \textbf{the} time, and we really don ' t have \textbf{the} immediate expertise at a technical level."In other words, we didn ' t have \textbf{the} technology to do what would be really interesting on that building.
And that would be to create a kind of three- dimensional bubble — a really interesting overcoat that would naturally ventilate, would breathe and would seriously reduce \textbf{the} energy loads.
Notwithstanding \textbf{the} fact that \textbf{the} building, as a green building, is very much a pioneering building.
And if I fast-forward in time, what is interesting is that \textbf{the} technology is now available and celebratory. \textbf{The} library of \textbf{the} Free University, which opened last year, is an example of that.
And again, \textbf{the} transition from one of \textbf{the} many thousands of sketches and \textbf{computer} images to \textbf{the} reality.
And a combination of \textbf{devices} here, \textbf{the} kind of heavy mass concrete of these book \textbf{stacks}, and \textbf{the} way in which that is enclosed by this skin, which enables \textbf{the} building to be ventilated, to consume dramatically less energy, and where it ' s really working with \textbf{the} forces of nature.
And what is interesting is that this is hugely popular by \textbf{the} people who use it.
Again, coming back to that thing about \textbf{the} lifestyle, and in a way, \textbf{the} ecological agenda is very much at one with \textbf{the} spirit.
So it ' s not a kind of sacrifice, quite \textbf{the} reverse.
I think it ' s a great — it ' s a celebration.
And you can measure \textbf{the} performance, in terms of energy \textbf{consumption}, of that building against a typical library.
If I show another aspect of that technology then, in a completely different context — this apartment building in \textbf{the} Alps in Switzerland.
Prefabricated from \textbf{the} most traditional of materials, but that material — because of \textbf{the} technology, \textbf{the} computing ability, \textbf{the} ability to prefabricate, make high-performance components out of timber — very much at \textbf{the} cutting edge.
And just to give a sort of glimpse of that technology, \textbf{the} ability to plot points in \textbf{the} sky and to transmit, to transfer that information now, directly into \textbf{the} factory.
So if you cross \textbf{the} border — just across \textbf{the} border — a small factory in Germany, and here you can see \textbf{the} guy with his \textbf{computer} screen, and those points in space are communicated.
And on \textbf{the} \textbf{left} are \textbf{the} cutting machines, which then, in \textbf{the} factory, enable those individual pieces to be fabricated and plus or minus very, very few millimeters, to be slotted together on site.
And then interestingly, that building to then be clad in \textbf{the} oldest technology, which is \textbf{the} kind of hand-cut shingles.
One quarter of a million of them applied by hand as \textbf{the} final finish.
And again, \textbf{the} way in which that works as a building, for those of us who can \textbf{enjoy} \textbf{the} spaces, to live and visit there.
If I made \textbf{the} leap into these new technologies, then how did we — what happened before that?
I mean, you know, what was life like before \textbf{the} mobile phone, \textbf{the} things that you take for granted?
Well, obviously \textbf{the} building still happened.
I mean, this is a glimpse of \textbf{the} interior of our Hong Kong bank of 1979, which opened in 1985, with \textbf{the} ability to be able to reflect sunlight deep into \textbf{the} heart of this space here.
And in \textbf{the} absence of \textbf{computers}, you have to physically model.
So for example, we would put models under an \textbf{artificial} sky.
For wind tunnels, we would literally put them in a wind tunnel and blast air, and \textbf{the} many kilometers of cable and so on.
And \textbf{the} turning point was probably, in our terms, when we had \textbf{the} first \textbf{computer}.
And that was at \textbf{the} time that we sought to redesign, reinvent \textbf{the} airport.
This is Terminal Four at Heathrow, typical of any terminal — big, heavy roof, blocking out \textbf{the} sunlight, lots of machinery, big pipes, whirring machinery.
And Stansted, \textbf{the} green alternative, which uses natural light, is a \textbf{friendly} place : you know where you are, you can relate to \textbf{the} outside.
And for a large part of its cycle, not needing electric light — electric light, which in turn creates more heat, which creates more cooling loads and so on.
And at that particular point in time, this was one of \textbf{the} few solitary \textbf{computers}.
And that ' s a little image of \textbf{the} tree of Stansted.
Not going back very far in time, 1990, that ' s our office.
And if you looked very closely, you ' d see that people were drawing with pencils, and they were pushing, you know, big rulers and \textbf{triangles}.
It ' s not that long ago, 17 years, and here we are now.
I mean, major transformation.
Going back in time, there was a lady called Valerie Larkin, and in 1987, she had all our information on one \textbf{disk}.
Now, every week, we have \textbf{the} equivalent of 84 million \textbf{disks}, which record our archival information on past, current and future projects.
That reaches 21 kilometers into \textbf{the} sky.
This is \textbf{the} view you would get, if you looked down on that.
But meanwhile, as you know, wonderful protagonists like Al Gore are noting \textbf{the} inexorable rise in temperature, set in \textbf{the} context of that, interestingly, those buildings which are celebratory and very, very relevant to this place.
Our Reichstag project, which has a very familiar agenda, I ' m sure, as a public place where we sought to, in a way, through a process of advocacy, reinterpret \textbf{the} relationship between society and politicians, public space.
And maybe its hidden agenda, an energy manifesto — something that would be free, completely free of fuel as we know it.
So it would be totally \textbf{renewable}.
And again, \textbf{the} humanistic sketch, \textbf{the} translation into \textbf{the} public space, but this very, very much a part of \textbf{the} \textbf{ecology}.
But here, not having to model it for real.
Obviously \textbf{the} wind tunnel had a place, but \textbf{the} ability now with \textbf{the} \textbf{computer} to explore, to plan, to see how that would work in terms of \textbf{the} forces of nature : natural ventilation, to be able to model \textbf{the} chamber below, and to look at biomass.
A combination of biomass, aquifers, burning vegetable \textbf{oil} — a process that, quite interestingly, was developed in Eastern Germany, at \textbf{the} time of its dependence on \textbf{the} Soviet Bloc.
So really, retranslating that technology and developing something which was so clean, it was virtually pollution-free.
You can measure it again.
You can compare how that building, in terms of its \textbf{emission} in tons of \textbf{carbon} dioxide per year — at \textbf{the} time that we took that project, over 7, 000 tons — what it would have been with natural gas and finally, with \textbf{the} vegetable \textbf{oil}, 450 tons.
I mean, a 94 percent reduction — virtually clean.
We can see \textbf{the} same processes at work in terms of \textbf{the} Commerce Bank — its dependence on natural ventilation, \textbf{the} way that you can model those gardens, \textbf{the} way they spiral around.
But again, very much about \textbf{the} lifestyle, \textbf{the} quality — something that would be more enjoyable as a place to work.
And again, we can measure \textbf{the} reduction in terms of energy \textbf{consumption}.
There is an \textbf{evolution} here between \textbf{the} projects, and Swiss Re again develops that a little bit further — \textbf{the} project in \textbf{the} city in London.
And this sequence shows \textbf{the} buildup of that model.
But what it shows first, which I think is quite interesting, is that here you see \textbf{the} circle, you see \textbf{the} public space around it.
What are \textbf{the} other ways of putting \textbf{the} same amount of space on \textbf{the} site?
If, for example, you seek to do a building which goes right to \textbf{the} edge of \textbf{the} pavement, it ' s \textbf{the} same amount of space.
And finally, you profile this, you cut grooves into it. \textbf{The} grooves become \textbf{the} kind of green lungs which give views, which give light, ventilation, make \textbf{the} building fresher.
And you enclose that with something that also is central to its appearance, which is a mesh of triangulated structures — again, in a long connection evocative of some of those works of Buckminster Fuller, and \textbf{the} way in which triangulation can increase performance and also give that building its sense of identity.
And here, if we look at a detail of \textbf{the} way that \textbf{the} building opens up and breathes into those atria, \textbf{the} way in which now, with a \textbf{computer}, we can model \textbf{the} forces, we can see \textbf{the} high pressure, \textbf{the} low pressure, \textbf{the} way in which \textbf{the} building behaves rather like an aircraft \textbf{wing}.
So it also has \textbf{the} ability, all \textbf{the} time, regardless of \textbf{the} direction of \textbf{the} wind, to be able to make \textbf{the} building fresh and efficient.
And unlike conventional buildings, \textbf{the} top of \textbf{the} building is celebratory.
It ' s a viewing place for people, not machinery.
And \textbf{the} base of \textbf{the} building is again about public space.
Comparing it with a typical building, what happens if we seek to use such design strategies in terms of really large-scale thinking?
And I ' m just going to give two images out of a kind of company research project.
It ' s been well known that \textbf{the} Dead Sea is dying. \textbf{The} level is dropping, rather like \textbf{the} Aral Sea.
And \textbf{the} Dead Sea is obviously much lower than \textbf{the} oceans and \textbf{seas} around it.
So there has been a project which rescues \textbf{the} Dead Sea by creating a pipeline, a pipe, sometimes above \textbf{the} \textbf{surface}, sometimes buried, that will redress that, and will feed from \textbf{the} Gulf of Aqaba into \textbf{the} Dead Sea.
And our translation of that, using a lot of \textbf{the} thinking built up over \textbf{the} 40 years, is to say, what if that, instead of being just a pipe, what if it is a lifeline?
What if it is \textbf{the} equivalent, depending on where you are, of \textbf{the} Grand Canal, in terms of tourists, habitation, desalination, \textbf{agriculture}?
In other words, water is \textbf{the} lifeblood.
And if you just go back to \textbf{the} previous image, and you look at this area of volatility and hostility, that a unifying design idea as a humanitarian gesture could have \textbf{the} affect of bringing all those warring factions together in a united cause, in terms of something that would be genuinely green and productive in \textbf{the} widest sense.
Infrastructure at that large scale is also inseparable from communication.
And whether that communication is \textbf{the} virtual world or it is \textbf{the} physical world, then it ' s absolutely central to society.
And how do we make more legible in this growing world, especially in some of \textbf{the} places that I ' m talking about — China, for example, which in \textbf{the} next ten years will create 400 new airports.
Now what form do they take?
How do you make them more \textbf{friendly} at that scale?
Hong Kong I refer to as a kind of analog experience in a digital age, because you always have a point of reference.
So what happens when we take that and you expand that further into \textbf{the} Chinese society?
And what is interesting is that that produces in a way perhaps \textbf{the} ultimate mega-building.
It is physically \textbf{the} largest project on \textbf{the} \textbf{planet} at \textbf{the} moment. 250 — excuse me, 50, 000 people working 24 hours, seven days.
Larger by 17 percent than every terminal put together at Heathrow — built — plus \textbf{the} new, un-built Terminal Five.
And \textbf{the} challenge here is a building that will be green, that is compact despite its size and is about \textbf{the} human experience of travel, is about \textbf{friendly}, is coming back to that starting point, is very, very much about \textbf{the} lifestyle.
And perhaps these, in \textbf{the} end, as celebratory spaces.
As Hubert was talking over lunch, as we sort of engaged in conversation, talked about this, talked about cities.
Hubert was saying, absolutely correctly,"These are \textbf{the} new cathedrals."And in a way, one aspect of this conversation was triggered on New Year ' s Eve, when I was talking about \textbf{the} Olympic agenda in China in terms of its green ambitions and aspirations.
And I was voicing \textbf{the} thought that — it just crossed my mind that New Year ' s Eve, a sort of symbolic turning point as we move from 2006 to 2007 — that maybe, you know, \textbf{the} future was \textbf{the} most powerful, innovative sort of nation. \textbf{The} way in which somebody like Kennedy inspirationally could say,"We put a man on \textbf{the} moon."You know, who is going to say that we cracked this thing of \textbf{the} dependence on \textbf{fossil} fuels, with all that being held to ransom by rogue regimes, and so on.
And that ' s a concerted platform.
It ' s more than one \textbf{device}, you know, it ' s \textbf{renewable}.
And I voiced \textbf{the} thought that maybe at \textbf{the} turn of \textbf{the} year, I thought that \textbf{the} inspiration was more likely to come from those other, larger countries out there — \textbf{the} Chinas, \textbf{the} Indias, \textbf{the} Asian-Pacific \textbf{tigers}.
Thank you very much.

% matched lemmas: agriculture, artificial, blue, carbon, computer, consumption, device, disk, ecology, emission, enjoy, evolution, fossil, friendly, left, object, occur, oil, planet, renewable, sea, selection, stack, sum, surface, sustainability, the, tiger, triangle, unclear, wing, yellow
\end{textsample}
